% Part: first-order-logic
% Chapter: completeness
% Section: outline

\documentclass[../../../include/open-logic-section]{subfiles}

\begin{document}

\iftag{FOL}
      {\olfileid{fol}{com}{out}}
      {\olfileid{pl}{com}{out}}

\olsection{证明概要}

完备性定理的证明有些复杂，初次阅读时很容易迷失方向。因此我们先概述该证明的思路。第一步是视角的转换，它使我们看到证明的途径。当我们将完备性理解为“每当 $\Gamma \Entails !A$ 时就有 $\Gamma \Proves !A$”时，甚至可能很难想出思路：因为要证明 $\Gamma \Proves !A$，我们必须找到一个推导，而假设 $\Gamma \Entails !A$ 看起来对此并无任何帮助。对于某些证明系统，或许可以直接构造出一个推导，但我们将采取一种略有不同的方法。所需的视角转换在于：完备性也可以表述为：“如果 $\Gamma$ 是一致的，那么它是可满足的。”也许我们可以利用 $\Gamma$ 中的信息，加上它是一致的这一假设，来构造一个满足 $\Gamma$ 中所有\iftag{FOL}{语句}{公式}的\iftag{FOL}{结构}{赋值}。毕竟，我们知道我们寻找的是何种\iftag{FOL}{结构}{赋值}：正是 $\Gamma$ 所描述的那种！

如果 $\Gamma$ 只包含\iftag{FOL}{原子语句}{命题变元}，那么为它构造一个模型是容易的。\iftag{FOL}{假设原子语句都具有 $\Atom{P}{a_1,\dots,a_n}$ 的形式，其中 $a_i$ 是常项符号。}{}我们需要做的只是给出\iftag{FOL}{%
  一个论域 $\Domain{M}$ 以及对 $P$ 的一个指派，使得 $\Sat{M}{\Atom{P}{a_1,\ldots,a_n}}$。但这并不难：
  令 $\Domain{M} = \Nat$，$\Assign{\Obj c_i}{M} = i$，并且对于每个 $\Atom{P}{a_1,\ldots,a_n} \in \Gamma$，将有序组 $\tuple{k_1,
    \dots, k_n}$ 放入 $\Assign{P}{M}$，其中 $k_i$ 是常项符号 $a_i$ 的指标（即 $a_i \ident \Obj c_{k_i}$）。}{%
  一个赋值 $\pAssign{v}$，使得对所有 $p \in \Gamma$ 都有 $\pSat{v}{p}$。只需令 $\pAssign{v}(p) = \True$ 当且仅当 $p \in \Gamma$。}

现在假设 $\Gamma$ 包含某个形如 $\lnot !B$ 的公式，其中 $!B$ 是原子公式。我们可能会担心，对 $\iftag{FOL}{\Struct{M}}{\pAssign{v}}$ 的构造会干扰使 $\lnot !B$ 为真的可能性。但此时 $\Gamma$ 的一致性便派上了用场：如果 $\lnot !B \in \Gamma$，那么 $!B \notin \Gamma$，否则 $\Gamma$ 就会不一致。而如果 $!B \notin \Gamma$，那么根据我们对 $\iftag{FOL}{\Struct{M}}{\pAssign{v}}$ 的构造，$\iftag{FOL}{\Sat/{M}}{\pSat/{v}}{!B}$，所以 $\iftag{FOL}{\Sat{M}}{\pSat{v}}{\lnot !B}$。至此并无问题。

如果 $\Gamma$ 包含复杂的非原子公式呢？比如说，它包含 $!A \land !B$。要使该公式为真，我们应当像 $!A$ 和 $!B$ 都属于 $\Gamma$ 那样继续构造。而如果 $!A \lor !B \in \Gamma$，那么我们必须使其中至少一个为真，即像其中之一属于 $\Gamma$ 那样继续构造。

这启示了如下思路：我们向 $\Gamma$ 中添加额外的公式，以使所得集合能够 保持一致； 确保对于每个可能的原子语句 $!A$，要么 $!A$，要么 $\lnot !A$ 属于该集合； 使得每当 $!A \land !B$ 属于该集合时，$!A$ 和 $!B$ 也均属于它，且当 $!A \lor !B$ 属于该集合时，$!A$ 或 $!B$ 中至少有一个也属于它，等等。我们持续进行这一过程（可能无限进行下去）。将如此添加的所有公式构成的集合称为 $\Gamma^*$。那么，我们上面的构造就会给出一个\iftag{FOL}{结构 $\Struct{M}$}{赋值 $\pAssign{v}$}，通过归纳法可证明它满足 $\Gamma^*$ 中的所有语句，从而由于 $\Gamma \subseteq \Gamma^*$，也满足 $\Gamma$ 中的所有语句。事实证明，保证和就足够了。一个满足的语句集合被称为\emph{完备的}。所以我们的任务将是把一致集合 $\Gamma$ 扩展成一个一致且完备的集合 $\Gamma^*$。

\iftag{FOL}{%
这个计划中有一个难点：如果 $\lexists[x][!A(x)] \in \Gamma$，我们本希望能选取某个常项符号~$c$ 并在此过程中加入 $!A(c)$。但我们怎么知道总能这样做呢？也许我们的语言中只有极少数常项符号，并且对它们中的每一个，都有 $\lnot !A(c) \in \Gamma$。此时就不能再加入 $!A(c)$ 了，因为那会使集合不一致，且我们也无从判断 $\Struct{M}$ 究竟应该使 $!A(c)$ 为真还是使 $\lnot !A(c)$ 为真。此外，$\Gamma$ 包含的语句也可能属于一个根本没有常项符号的语言（例如集合论的语言）。

这个问题的解决方案是，一开始就直接加入无穷多个常项，并添加一些语句，以正确的方式将这些常项与量词联系起来。（当然，我们必须验证这样做不会引入不一致性。）

如果我们仅仅在原子语句中用到常项符号，那么原先的构造是可行的。但语言也可能包含函数符号。那样的话，要找到合适的、定义在~$\Nat$ 上的函数来指派给这些函数符号，以确保一切按预期运作，可能会很棘手。所以，这里还有另一个技巧：不再用 $i$ 来解释 $\Obj c_i$，而是直接把常项符号的集合当作论域。这样一来，$\Struct M$ 就可以把每个常项符号指派给它自身：$\Assign{\Obj c_i}{M} = \Obj c_i$。但为何不更彻底一些：让 $\Domain{M}$ 是语言中所有的\emph{项}呢！如果我们这样做，那么为函数符号指派函数（这些函数以项为自变元，值也是项）就有了一个显而易见的方式：我们指派给函数符号~$\Obj f^n_i$ 这样一个函数，它以 $n$ 个项 $t_1$，……，$t_n$ 为输入，输出项 $\Obj f^n_i(t_1, \dots, t_n)$ 作为值。

最后剩下的难题是如何处理~$\eq$。谓词符号~$\eq$ 有一个固定的解释：$\Sat{M}{\eq[t][t']}$ 当且仅当 $\Value{t}{M} = \Value{t'}{M}$。现在，如果我们让一个项~$t$ 的值就是 $t$ 本身，那么这种结构将使得形如 $\eq[t][t']$ 的语句\emph{无一}为真，除非 $t$ 和 $t'$ 是同一个项。这显然是个问题，因为基本上在带函数符号的语言中，任何有趣的理论都会把某些 $t$ 和 $t'$ 不是相同项的 $\eq[t][t']$ 当作定理（例如，在算术理论中：$\eq[(\Obj 0+ \Obj 0)][\Obj 0]$）。为了解决这个问题，我们修改~$\Struct M$ 的论域：不再用单个的项作为~$\Domain{M}$ 中的对象，而是用项的集合，并且每个集合都包含所有那些被~$\Gamma$ 中的语句要求相等的项。例如，如果 $\Gamma$ 是一个算术理论，那么这些集合中就会有一个包含：$\Obj 0$、$(\Obj 0 + \Obj 0)$、$(\Obj 0 \times \Obj 0)$，等等。这个集合就是我们将指派给 $\Obj 0$ 的集合，并且可以证明，这个集合同时也是其中所有项（例如也包括 $(\Obj 0 + \Obj 0)$）的值。因此，语句 $\eq[(\Obj 0+ \Obj 0)][\Obj 0]$ 在这个修改后的结构中将为真。}{}

那么，我们要做的是：首先研究完备一致集的性质，
特别地，我们将证明，一个完备一致集包含 $!A \land !B$ 当且仅当它同时包含
$!A$ 和~$!B$，包含 $!A \lor !B$ 当且仅当它包含二者中至少一个，等等
（\olref[ccs]{prop:ccs}）。\iftag{FOL}{接着我们定义并研究语句的“饱和”集。
一个饱和集包含连接每个量化语句与其实例的条件句
（\olref[hen]{defn:henkin-exp}）。
我们证明任何一致集~$\Gamma$ 总可以扩充为一个饱和集~$\Gamma'$
（\olref[hen]{lem:henkin}）。
如果一个集合是一致的、饱和的并且是完备的，那么它还具有如下性质：
它包含 \iftag{prvEx}{$\lexists[x][!A(x)]$ 当且仅当它包含某个闭项~$t$ 的 $!A(t)$}{}\iftag{defEx,defAll}{}{ 并且 }
\iftag{prvAll}{$\lforall[x][!A(x)]$ 当且仅当它对所有闭项~$t$ 都包含 $!A(t)$}{}
（\olref[hen]{prop:saturated-instances}）。}{}
然后，我们取那个\iftag{FOL}{ 饱和}{}一致集~$\iftag{FOL}{\Gamma'}{\Gamma}$，并证明它可以扩充为
一个\iftag{FOL}{ 饱和、一致且}{ 一致且}完备的集~$\Gamma^*$
（\olref[lin]{lem:lindenbaum}）。
这个集合 $\Gamma^*$ 将用来定义我们的
\iftag{FOL}{项模型~$\Struct M(\Gamma^*)$}{赋值~$\pAssign v(\Gamma^*)$}。
\iftag{FOL}{项模型以所有闭项的集合作为其论域，其谓词符号的解释由
$\Gamma^*$ 中的原子语句给出}{赋值由 $\Gamma^*$ 中的命题变元确定}
（\olref[mod]{defn:termmodel}）。
我们将利用\iftag{FOL}{ 饱和、}{}完备一致集的性质来证明，
确实有$\iftag{FOL}{\Sat{M(\Gamma^*)}}{\pSat{v(\Gamma^*)}}{!A}$ 当且仅当 $!A \in \Gamma^*$（\olref[mod]{lem:truth}），
因此特别地，$\iftag{FOL}{\Sat{M(\Gamma^*)}}{\pSat{v(\Gamma^*)}}{\Gamma}$。
\iftag{FOL}{最后，我们将考虑如果 $\Gamma$ 还包含~$\eq$，应如何定义项模型
（\olref[ide]{defn:term-model-factor}），并证明它满足~$\Gamma^*$
（\olref[ide]{lem:truth}）。}{}

\end{document}